\documentclass[11pt]{article}
\usepackage[utf8]{inputenc}
\usepackage{amsmath,amssymb}
\usepackage{hyperref}
\title{Efficient Extreme Weather Attribution: A Resource-Aware Approach}
\author{Anonymous}
\date{}
\begin{document}
\maketitle

\begin{abstract}
This paper studies Extreme Weather Attribution. Grounded in a real, citable literature \cite{ref1,ref2,ref3,ref4,ref5,ref6,ref7,ref8},
we propose a focused method and report \textbf{illustrative} experiments.
All references below are real and verifiable (OpenAlex/arXiv); the reported
numbers are simulated to illustrate intended behaviour.
\end{abstract}

\section{Introduction}
Extreme Weather Attribution is an active research area. This work targets a concrete gap identified from
the recent literature and contributes a method together with an evaluation protocol.

\section{Related Work (real literature)}
The following works, retrieved from public bibliographic APIs, frame our study:
\begin{itemize}
\item Epanechnikov (1969) — "Non-Parametric Estimation of a Multivariate Probability Density", Theory of Probability and Its Applications (cited 1838).
\item Fischer et al. (2015) — "Anthropogenic contribution to global occurrence of heavy-precipitation and high-temperature extremes", Nature Climate Change (cited 1418).
\item Trenberth et al. (2015) — "Attribution of climate extreme events", Nature Climate Change (cited 974).
\item Guthrie (1993) — "Faces <i>in the</i> Clouds" (cited 905).
\item Pall et al. (2011) — "Anthropogenic greenhouse gas contribution to flood risk in England and Wales in autumn 2000", Nature (cited 892).
\item Ummenhofer et al. (2017) — "Extreme weather and climate events with ecological relevance: a review", Philosophical Transactions of the Royal Society B Biological Sciences (cited 886).
\end{itemize}

\section{Method}
We reduce the dominant computational cost via adaptive resource allocation, activating only the components needed per input. We instantiate this idea for Extreme Weather Attribution and analyse its cost/quality trade-off.

\section{Experiments (Illustrative)}
\begin{center}
\begin{tabular}{lcc}
System & Primary Metric & Efficiency \\ \hline
Strong baseline & 0.812 & 1.00$\times$ \\
Ablation & 0.828 & 0.92$\times$ \\
\textbf{Ours} & \textbf{0.851} & \textbf{0.71$\times$}
\end{tabular}
\end{center}
\emph{Metrics simulated to illustrate the intended effect; not measured.}

\section{Conclusion}
We connected a real literature to a concrete method for Extreme Weather Attribution. Future work will
validate the illustrative results with real experiments.

\begin{thebibliography}{9}
\bibitem{ref1} V. A. Epanechnikov. Non-Parametric Estimation of a Multivariate Probability Density. \emph{Theory of Probability and Its Applications}, 1969. DOI:10.1137/1114019
\bibitem{ref2} Erich Fischer, Reto Knutti. Anthropogenic contribution to global occurrence of heavy-precipitation and high-temperature extremes. \emph{Nature Climate Change}, 2015. DOI:10.1038/nclimate2617
\bibitem{ref3} Kevin E. Trenberth, John Fasullo, Theodore G. Shepherd. Attribution of climate extreme events. \emph{Nature Climate Change}, 2015. DOI:10.1038/nclimate2657
\bibitem{ref4} Stewart Guthrie. Faces <i>in the</i> Clouds. \emph{}, 1993. DOI:10.1093/oso/9780195069013.001.0001
\bibitem{ref5} Pardeep Pall, T. Aina, Dáithí A. Stone et al.. Anthropogenic greenhouse gas contribution to flood risk in England and Wales in autumn 2000. \emph{Nature}, 2011. DOI:10.1038/nature09762
\bibitem{ref6} Caroline C. Ummenhofer, Gerald A. Meehl. Extreme weather and climate events with ecological relevance: a review. \emph{Philosophical Transactions of the Royal Society B Biological Sciences}, 2017. DOI:10.1098/rstb.2016.0135
\bibitem{ref7} Peter A. Stott, Nikolaos Christidis, Friederike E. L. Otto et al.. Attribution of extreme weather and climate‐related events. \emph{Wiley Interdisciplinary Reviews Climate Change}, 2015. DOI:10.1002/wcc.380
\bibitem{ref8} Nikolaos Christidis, Gareth S. Jones, Peter A. Stott. Dramatically increasing chance of extremely hot summers since the 2003 European heatwave. \emph{Nature Climate Change}, 2014. DOI:10.1038/nclimate2468
\end{thebibliography}
\end{document}
